% !TeX program = xelatex
% ============================================================================
%  第 12 课　课堂单页（学生用，单独打印一张 A4）
%  编译：XeLaTeX，两遍即可
%  用途：把费马小定理的证明，自己一步一步写出来
% ============================================================================

\documentclass[11pt]{ctexart}
\usepackage[a4paper,left=1.7cm,right=1.7cm,top=1.4cm,bottom=1.4cm]{geometry}
\usepackage{amsmath,amssymb}
\usepackage{booktabs}
\usepackage{array}
\usepackage[dvipsnames]{xcolor}

\setCJKmainfont{SimSun}[BoldFont=SimHei,ItalicFont=KaiTi]
\IfFontExistsTF{TeX Gyre Pagella}{\setmainfont{TeX Gyre Pagella}}{\setmainfont{Latin Modern Roman}}

\definecolor{pagegreen}{RGB}{34,139,34}
\renewcommand{\arraystretch}{1.4}
\setlength{\parindent}{0pt}
\pagestyle{empty}

\newcount\xhcnt
\newcommand{\liukong}[1]{%
  \par\vspace{4pt}%
  \noindent\begin{minipage}{\linewidth}%
    \setlength{\parindent}{0pt}\setlength{\parskip}{0pt}%
    \xhcnt=#1\relax
    \loop\ifnum\xhcnt>0
      \noindent\textcolor{black!35}{\rule{\linewidth}{0.5pt}}\par\vspace{0.78cm}%
      \advance\xhcnt by -1\relax
    \repeat
    \vspace{-0.78cm}%
  \end{minipage}%
  \par\vspace{3pt}%
}

\newcommand{\biaoti}[2]{%
  {\color{pagegreen}\rule{\textwidth}{1.2pt}}\\[6pt]
  {\heiti\Large 第 #1 课\quad #2}\hfill{\small\songti 课堂单页}\\[6pt]
  {\color{pagegreen}\rule{\textwidth}{0.5pt}}\\[4pt]
  {\small 姓名\underline{\hspace{3.2cm}}\qquad 班级\underline{\hspace{2.4cm}}\qquad 座号\underline{\hspace{1.6cm}}}
  \vspace{0.3em}
}

\begin{document}

\biaoti{12}{费马小定理的证明}

\section*{一、取 \(p=7\)，\(a=3\)，先把表填满}

\begin{center}
{\small
\begin{tabular}{@{}lcccccc@{}}
  \toprule
  原来的数 & \(1\) & \(2\) & \(3\) & \(4\) & \(5\) & \(6\) \\
  \midrule
  乘 \(3\) 之后 & \(3\) & \(6\) & \(9\) & \(12\) & \(15\) & \(18\) \\
  \rule{0pt}{2.8ex}取余数 & & & & & & \\
  \bottomrule
\end{tabular}}
\end{center}

\noindent 最后一行是 \(1,2,3,4,5,6\) 的一次\underline{\hspace{4em}}。

\section*{二、把两行各自乘起来}

\[
  1\times2\times3\times4\times5\times6=\underline{\hspace{2.2em}},
  \qquad
  3\times6\times2\times5\times1\times4=\underline{\hspace{2.2em}} .
\]

\noindent 两个乘积一模一样，因为下排只是上排的重排。

\section*{三、换一个角度乘，再划掉}

\[
  3\times6\times9\times12\times15\times18
  =\underline{\hspace{2.2em}}\times 6! ,
\]

\noindent 所以两边取模 \(7\)：

\[
  \underline{\hspace{2.2em}}\times6!\equiv 6! \pmod 7 .
\]

\noindent 两边都有的 \(6!\) 能不能划掉？\underline{\hspace{2.6em}}\qquad
为什么？\underline{\hspace{6em}}

\vspace{0.2em}
\noindent 划完得到：\(3^{6}\equiv\underline{\hspace{2.6em}}\pmod 7\)。

\section*{四、换成字母，把定理写出来}

\begin{center}
{\small
\begin{tabular}{@{}ll@{}}
  \toprule
  \textbf{具体}（\(p=7\)，\(a=3\)） & \textbf{一般} \\
  \midrule
  \(3,\ 6,\ 9,\ \dots,\ 18\) & \underline{\hspace{5.6em}} \\
  余数是 \(1\text{—}6\) 的重排 & \underline{\hspace{5.6em}} \\
  \(3^{6}\times6!\equiv6!\) & \underline{\hspace{5.6em}} \\
  \(3^{6}\equiv1\pmod 7\) & \underline{\hspace{5.6em}} \\
  \bottomrule
\end{tabular}}
\end{center}

\vspace{0.3em}
\noindent{\heiti 费马小定理：}设 \(p\) 是素数，\(a\) 不被 \(p\) 整除，则
\[
  a^{\,p-1}\equiv\underline{\hspace{2.2em}} \pmod p .
\]

\end{document}